% !PRINTED 380
\weberpage{380}
%<<CONT

A complex number $\omega = \xi + i\eta$ with positive $\eta$ is called \gesperrt{reduced} if
\weq{2}{-\tfrac12 \leqq \xi \leqq \tfrac12, \quad \xi^2+\eta^2 \geqq 1.}
If one regards $\xi$, $\eta$ as rectangular coordinates in a plane, then every number $\omega$ is represented by a point of this plane, and the position of the points that correspond to the reduced numbers is represented by the region shaded in Fig. 26 (with inclusion of the boundaries), which we call the fundamental region.

% SOURCE FIGURE CROP: Bd. I registered scan 406 (printed p. 380), 600 dpi,
% crop x=350,y=1005,w=1180,h=900; source-faithful raster, Fig. 26.
\weberfigureimage{weber-b1-fig26.png}{0.48}

From the conditions (2) it then further follows, as is also easily confirmed in the figure, that
\weq{3}{\eta^2 \geqq \tfrac34.}
If one now passes from a number $\omega$ to an equivalent number $\omega_1$, formula (16) of the preceding section gives
\weq{4}{\omega_1 = \frac{(\alpha\xi+\beta)+i\alpha\eta}{(\gamma\xi+\delta)+i\gamma\eta} = \frac{[(\alpha\xi+\beta)(\gamma\xi+\delta)+\alpha\gamma\eta^2]+i\varepsilon\eta}{(\gamma\xi+\delta)^2+\gamma^2\eta^2},}
and this formula shows that $\eta_1$ has the same sign as $\eta$ or the opposite one, according as $\varepsilon = +1$ or $=-1$, that is, according as the equivalence is a proper or an improper one.

If we restrict ourselves, then, to numbers with positive imaginary part, only proper equivalence comes into consideration, and we now prove the fundamental theorem:

1. \gesperrt{Every imaginary quadratic irrational number $x+y\sqrt{d}$ with positive imaginary part is equivalent to a reduced number.}

If we understand by $\alpha$ the integer nearest to $\xi$, then $\omega-\alpha$ will satisfy the first of the conditions (2), namely that $-\tfrac12 \leqq \xi-\alpha \leqq \tfrac12$.

If now
\weq{5}{r^2 = (\xi-\alpha)^2+\eta^2 = (\omega-\alpha)(\omega'-\alpha)}
is greater than or equal to 1, then $\omega-\alpha$, which is equivalent to $\omega$, is already reduced; otherwise we set, applying a continued-fraction expansion,
%CONT>>
